\chapter{Ocean Vuong 談寫作的藝術}

本章追隨的是 Ocean Vuong 與主持人之間的一場對談，收錄於 LazyLearn 的寫作與說話系列中，由 LazyingArt LLC 策劃。講座一開始，便把讚嘆轉化成方法：驚異、新鮮感、著迷、神奇，這一切都被轉譯成一個實作性的工藝問題。我們要如何讓語言重新看見？這場講座沒有任何可經驗證的黑板或投影片數學留存，因此下文中的顯示公式都只是根據逐字稿壓縮出來的口頭結構，而不是可見畫面的逐字轉錄。它們之所以在這裡有用，是因為整場講座本身反覆以對照、轉換、門檻與因果鏈來思考。

\section{隱喻、觀察，以及與摹擬的第一次決裂}

主持人一開始，先點出 Vuong 作品中的一種品質，然後才開始問方法。他說，文字裡有驚異，但這種驚異不只存在於文字本身，也存在於觀看的方式裡；於是他問：這種觀看是否能像一塊肌肉一樣被鍛鍊？Vuong 的第一個回答立刻就進入技術層面：隱喻就是這個問題變得可以教的地方。如果學生問該怎麼寫出好的隱喻，那第一個答案不是靈光，也不是裝飾，而是觀察。你要長久地看著世界。至於其他東西，他說，都是安排與句法。

\begin{definition}
令 tenor 為被描寫之物，vehicle 為把描寫帶往他處之物，那麼 Vuong 對隱喻的開場示意便可以寫成
\begin{equation}
(\text{tenor},\text{vehicle}) \Longrightarrow \text{新的對應} \Longrightarrow \text{重新看見}.
\end{equation}
一個隱喻之所以成功，是因為它所建立的對應會迫使讀者重新知覺 tenor，而不是僅僅把它辨認出來。
\end{definition}

也正因如此，講座立刻求助於亞里斯多德對兩種句子功能的區分：
\begin{align}
\text{摹擬} &:\ \text{已知場景} \mapsto \text{可辨認的描寫},\\
\text{創造} &:\ \text{已知場景} \mapsto \text{被陌生化的過程}.
\end{align}
摹擬式句子讓我們回到早已知道如何命名的事物；詩性的句子則打斷這種自動命名。依這個精神，Vuong 的開場規則可以寫成
\begin{equation}
\text{隱喻} \neq \text{裝飾}, \qquad \text{隱喻} = \text{觀察} + \text{安排／句法}.
\end{equation}

講座裡第一個完整例子，是 Isaac Babel 在 \textit{Red Cavalry} 裡寫落日的句子。Vuong 把一條報紙式的句子——「山丘上一輪紅色的黃昏落日」——和 Babel 那輪「彷彿被斬首般在山間滾動的低紅太陽」對照起來。重點不只是說 Babel 很鮮明，而是：歷史情境本身被帶進了這個句子。我們不需要外加註解告訴我們 Babel 是從戰爭中寫作，因為影像本身就完成了這件事。它同時也改變了場景被感受到的運動方式。

\paragraph{Worked Example: Babel's Sunset.}
\begin{align}
\text{落日} + \text{斬首意象} &\Longrightarrow \text{戰爭語境被嵌入知覺之中},\\
\text{被嵌入的暴力} &\Longrightarrow \text{落日速度被改變}.
\end{align}
\begin{enumerate}
\item 我們先從一個完全可辨認的場景開始：山丘上的落日。
\item 接著加入一個不只是裝飾落日，而是把暴力帶進落日中的意象。
\item 這個加入的意象，於是把歷史意識直接嵌進知覺本身。
\item 那輪落日現在不只是落下，而像是在滾動、墜落、失去頭顱；第二個分句改變了整個場景的速度。
\item 因此，這句話做的事情不只是描寫，而是重新組織讀者的觀看。
\end{enumerate}

Vuong 為這種企圖給出了一句故意嚴苛的口號：
\begin{equation}
\text{稱職的句子} \neq \text{物種以前從未擁有過的句子}.
\end{equation}
左邊那一項並不是無用的，我們當然仍然需要支架式句子；但講座從一開始就堅持，寫作者更深的任務，不只是製造稱職的辨認，而是讓一個新的知覺事件得以發生。

\subsection{Question \& Answer}

我們如何寫出一個不只是模仿世界的隱喻？

不是先去追逐機智。我們先看，看到那個物件原先現成的名稱開始顯得不夠用。接著，我們去尋找一個來自持久觀察的 vehicle，而不是一個現成、被大量使用的相似物。整個序列是非常實作性的：
\begin{enumerate}
\item 先辨認出 tenor，並持續觀看它，直到它不再顯得泛泛。
\item 從持久觀察裡選出 vehicle，而不是從套裝相似性裡選。
\item 建立一條會改變場景、而不只是潤飾場景的對應。
\item 只有當讀者真的被迫看見，而不是只是認出來時，這一行才值得被保留。
\end{enumerate}
這也就是為什麼一個強而有力的隱喻，往往可能需要好幾年才到來。真正困難的，不只是句法，而是你是否掙得了那條對應。

\section{工作坊作為辨識，而不是糾正}

當講座已經示範了活的句子能做什麼之後，主持人便把鏡頭拉近。如果寫作可能膽怯，也可能活著，那我們該如何教作者看見自己作品中的差別？Vuong 的回答，是用「工作坊作為辨識」取代「工作坊作為糾正」。這個轉換很細，卻是決定性的。工作坊並不會只因為產生更多評論、更多草稿、更多程序，就自動變得有幫助。

講座把這種對照壓縮成
\begin{align}
\text{模式辨識} &\Longrightarrow \text{自我辨識} \Longrightarrow \text{更聰明的修訂},\\
\text{教條式糾正} &\Longrightarrow \text{作品被同質化或被摧毀}.
\end{align}
這一點之所以重要，是因為更廣泛的文化很容易讓我們用工廠的方式思考：把作品送進機器，加上更多流程，它就應該會變得更好。Vuong 拒絕這個幻想。有些工作坊會摧毀作品，有些草稿甚至超過了作品原本的高點。看似生產力的東西，其實可能只是瓦礫。

講座最清楚的一條形式性定義也就在這裡出現：
\begin{equation}
\text{句子} = \text{經由句法過濾的意識}.
\end{equation}
如果這句話成立，那麼工作坊就必須先問：這份草稿正在顯露的是哪一種意識？哪些意象一再回來？哪些動詞總是落在線行之間？哪一段突然換了時態？哪些介系詞總是把句子向前彈出去？Vuong 在課堂裡給的例子之所以刻意具體，正是因為他的重點不是抽象的「主題」，而是那顆心智此刻正在如何形式性地運動。

這也就是為什麼講座會再次回到「過度生產」的危險。一位詩人若一年裡每天都寫一首詩，未必會得到更多可用作品；他也可能只會得到一片比重新開始更難清理的廢墟。辨識則問的是另一個更簡單的問題：這裡其實已經在發生什麼？一旦這個問題變得嚴肅，修訂就不再是套用教條，而是澄清傾向。

那位日本植物學家的故事，正把這件事轉化成一種方法。有人問他怎麼找到那麼多藥草，他說自己並不是去找那些看起來像現成藥材的植物；他只是去找那些對自己來說陌生的新植物，然後希望其中某些可能有藥性。有些當然不是，有些甚至有毒。但真正擴大整個場域的，不是像既定模型，而是對新東西的追尋。Vuong 很明顯是把這當成寫作課的一種模型提出來的。

\subsection{Question \& Answer}

為什麼更多回饋與更多起草，反而可能讓一篇作品變得更糟？

因為更多流程不等於更好的辨識。當修訂被作品本身未曾掙得的既有規則支配時，它就會失敗。Vuong 的替代方案雖然也有程序，卻不是機械式的：
\begin{enumerate}
\item 在提出修補方案之前，先讀出模式。
\item 先命名那些反覆出現的形式傾向：措辭、時態、句法、意象與行行之間的運動。
\item 追問作品究竟在哪些地方突然意外地活了起來。
\item 修訂要朝向這些發現，而不是朝向一般化的「正確」。
\end{enumerate}
在這個模型裡，工作坊不是用來製造詩的，而是幫助作者更快抵達詩原本已經試圖成為的東西。

\section{句子是如何被馴化的}

中斷之後，講座並不是靠簡單回顧來接續，而是靠張力來接續。主持人問起新穎、驚奇、新鮮感，以及「追求品質」本身會不會反過來變成一種限制。Vuong 的回答則是把畫面拉大。他說，現代句子並不是一條永恆的好寫作標準，而是一種歷史產物。

而這段更大的歷史，始於一個更前面的問題：我們究竟說「文學」是在說什麼？Vuong 堅持，這個範疇本身其實是晚近的。在現代學院與其制度分類出現之前，詩、故事、書信、演說與實用文字之間的邊界，都要鬆動得多。就連小說本身，也不是一直都被當成某種嚴肅文學工具。它後來在美國語境中，特別是在內戰之後，被賦予某種道德高度，而這又恰好與另一件事同時發生：新聞寫作的標準化。

這一段歷史對照，可以示意地寫成：
\begin{align}
\text{演說式／維多利亞句子} &\Longrightarrow \text{延宕} + \text{從屬} + \text{堆積},\\
\text{報紙句子} &\Longrightarrow \text{簡短} + \text{清晰} + \text{效率} + \text{標準化}.
\end{align}
前者屬於講道、演說、公共論辯，也屬於一種仰賴節奏與拖延來維持注意力的文化。後者則屬於資訊傳遞，它必須可靠、清楚、快速，並且足夠簡短，才能與發行與廣告共存。

也正是因此，講座裡提到 Churchill 的首語重複，不是無關緊要的旁支，而是對舊式句子的現場記憶。``We shall fight...'' 透過延後支付，在公共語勢中逐步堆積力量。接著，Vuong 又審慎而清楚地回到逐字稿中那段雖有破損、但意思穩定的部分：二十世紀散文吸收了報紙模型，而像 Hemingway、Crane、London、Orwell 這樣的作家，逐漸成為我們文化中所謂「好散文」的那種簡練、高效句子的代表。

當主持人提出「直角」這個比喻時，論點便開始變得更尖。今天很多散文——他說——都像尺畫出來的一樣筆直、方正、粗糙，卻又過度修整。Vuong 立刻接住這個類比：
\begin{equation}
\text{工業式直角} \Longrightarrow \text{文化對標準化形式的偏好}.
\end{equation}
這當然不是幾何學，而是一種描述：工業現代性如何教我們去偏愛乾淨邊角、可重複性，以及可大規模複製的表面。散文也因此受到同樣的訓練。句子開始追求「隱形」、高效，並對作者性的壓力感到懷疑。

從這裡，講座自然走向當代例子。Microsoft Word 會用修正標記蓋滿 Shakespeare。AI 的出現，並不是人類語言忽然遭遇背叛，而是某個更古老的同質化計畫的延續。早在自動化散文之前，我們就已經在訓練句子朝向彼此相似的聲音。

\subsection{Question \& Answer}

為什麼當代散文常常聽起來被標準化了，即使我們嘴上一直在讚美創新？

因為我們的制度會在事後讚美創新，卻又在程序上壓抑創新。我們一方面教授那些大膽的經典作者，一方面卻又要求仍在寫作的學生聲音不要太顯眼、太特異、太不像既有商品。編輯、評審文化、軟體與市場恐懼，全都獎勵「可辨認性」。結果不是散文的死亡，而是散文被允許發聲的範圍被縮窄了。

\section{陌生化、門檻，以及重新學會觀看}

一旦這種窄化被診斷出來，主持人便直接點名 disenchantment。Vuong 的回答則引向中段講座的那個反命題：estrangement。此時，Viktor Shklovsky 變得居於核心。這裡最強的一句話是：事物本身根本沒有什麼「陳腔濫調」可言。陳腔濫調只存在於被耗盡了的處理方式裡。如果我們開始把玫瑰、祖母、廚房、山、花朵全都當成「已經被寫過了」的禁物，那我們最後會把整個世界都逐出文學。

最基本的結構便是：
\begin{equation}
\text{熟悉的物件} + \text{意外的安置} \Longrightarrow \text{新鮮的知覺}.
\end{equation}
也因此，講座中最尖銳的例子才會是：
\begin{equation}
\text{伴娘髮間的玫瑰} \neq \text{Mike Tyson 耳上的玫瑰}.
\end{equation}
玫瑰本身完全沒變，改變的是處理方式。所以這裡的教訓不是「不要再寫玫瑰」，而是「重新觀看那朵玫瑰」。

Shklovsky 引 Tolstoy 的段落，則把這一點說得更深。Tolstoy 記不清自己是否已經撢過沙發上的灰，因為那個動作太例行、太無意識。若生活只以例行方式被度過，那它便彷彿根本沒有被活過。文學之所以在這裡有其重要性，正是因為它打斷例行。它恢復了「看見」與「僅僅認出來」之間的差別。

主持人在這部分給出的那些類比尤其重要，因為它們讓抽象保持著具體。Bierstadt 讓山再次變得可見。Monet 的睡蓮與 Van Gogh 的花，不只是展示物件本身，也展示其中的能量。接著，主持人又把話題帶到影片剪輯與早期攝影：當我們把鏡頭放大、把格數拆細，原本看似單一而平滑的動作，就會顯出由許多不同瞬間構成。這正是 Vuong 對知覺的要求：寫作必須把世界放慢到足以讓我們遇見那些中間狀態。

於是，講座再次回到亞里斯多德：
\begin{equation}
\text{花苞} \leadsto \text{玫瑰}, \qquad \text{創造} = \text{命名狀態之間的門檻區域}.
\end{equation}
玫瑰是一個被命名的東西，花苞也是一個被命名的東西；但在兩者之間，其實存在無數個過渡階段。這些階段並不會因為難以命名就變得不真實。恰恰相反，Vuong 說，wonder 就住在那裡。Heidegger 的門檻語言之所以有幫助，是因為它迫使我們問一個日常命名總是急著跳過的問題：玫瑰究竟是在何處，才真正變成玫瑰的？

這也使得 Vuong 對亞里斯多德的援引，比教科書式的詩學說明更強。他不是拿 Aristotle 來當歷史擺設，而是拿他來標記那個正確的地方：就在摹擬性命名失效、而活生生過程開始的位置。

\subsection{Question \& Answer}

一個題材本身會陳腔濫調嗎？還是只有當我們未能把它陌生化時，它才會死掉？

在這場講座裡，題材本身沒有罪，真正有罪的是處理方式。某個題材之所以會死掉，只是因為我們用陳舊的辨認方式去接近它。實作上的規則因此是：
\begin{enumerate}
\item 不要把熟悉物件禁掉。
\item 拒絕你對它最先想到的那種處理方式。
\item 換一個位置、換一個時間、或換一種命名。
\item 留在那裡，直到單純的辨認再次打開成真正的看見。
\end{enumerate}
陌生化不是拒絕世界，而是把世界從習慣裡救回來。

\section{句法、餘響，以及閱讀的模式邏輯}

一旦知覺重新被打開，講座便轉向新的難題。現在已經不夠只問「我們怎麼看」，我們還得問：這種看見要如何在讀者身上存活下去？Vuong 很直接地把問題說出來：他不那麼關心「如何抓住一個讀者」，他更關心的是「如何與一個讀者待在一起」。主持人則不斷把這個意思翻譯成讀者經驗中的語言，於是產生了整場講座中最清楚的幾條說法。

Vuong 的比例當然是粗略的，但它很有用：
\begin{equation}
80\% \text{ 的觀看與思考}, \qquad 20\% \text{ 的句法}.
\end{equation}
他立刻修正了聽者可能會產生的誤解。最後那 20\% 才是「一切」，因為句法就是下載機制。因此，講座中的工作模型便是
\begin{equation}
\text{知覺} + \text{句法塑形} \Longrightarrow \text{讀者心中的痕印／餘響}.
\end{equation}

Siken 的例子，正好展示了這種機制如何在實際寫作中運作。群星變成了 ``little boats rowed out too far.'' 其 tenor 其實很熟悉：
\begin{align}
\text{tenor} &= \text{星星},\\
\text{vehicle} &= \text{划得太遠的小船}.
\end{align}
但這條對應徹底改變了整個象徵系統。原本宏大、神話性的星星，變成了親密、孤單、太晚、脆弱的東西。這一句話做的，不只是提供一個新鮮的比較，而是改變了夜空的情感尺度。

Ben Lerner 接著又提供了同一種要求的一個更殘酷的教學版本。學生寫出一句「還不錯」的句子，而 Google 一查，卻發現已經有幾十萬人寫過它。這個打擊之所以讓人難忘，正因為它把講座一開始的標準磨得更尖：若一條句子讓知覺停留在原地，那麼稱職還遠遠不夠。

從這裡，講座自然走向 haunting。Browning 的 \textit{Meeting at Night} 之所以一直留在 Vuong 心裡，正是因為它無法被壓平為一個簡單的概念。Good Will Hunting 裡公園長椅那一場戲，也是一樣。若只是平板地講一句關於愛與經驗的話，概念仍然可以被傳達；但真正有力的，是那個場景如何進入一種脆弱、帶著圖像感的句法之中。主持人在這一段扮演的角色格外重要：他不斷把剛才發生的事，用更平白的話重新說一次。前一刻，那還只是辨認；下一刻，它就已經變成一種看法。

接著，整個論證開始從句子擴大到藝術本身：
\begin{equation}
\text{模式} \Longrightarrow \text{不是滿足期待，就是拒絕期待}.
\end{equation}
因為句子是一種線性技術，它在時間中展開；音樂如此、剪輯如此、DJ set 如此、滑板影片也是如此。一旦有序列，就有期待；一旦有期待，創作者便獲得了兩種基本資源：
\begin{equation}
\text{滿足／拒絕期待} \Longrightarrow \text{愉悅} + \text{驚奇} + \text{陌生化}.
\end{equation}

\paragraph{Pattern Rule.}
\begin{enumerate}
\item 線性藝術透過在時間中展開，生成期待。
\item 創作者可以滿足那種期待，也可以暫時扣住它。
\item 延遲會增加注意力，因為觀者能感到某種完成被暫時保留了。
\item 釋放，或對釋放的拒絕，會創造愉悅與驚奇。
\item 到了這一刻，我們不再只是追蹤內容，而是開始感受模式本身的智能。
\end{enumerate}

這也就是為什麼 Vuong 可以如此自然地在文學、街頭籃球的假動作、滑板片段，以及 DJ 的 timing 之間來回穿梭。這些類比不是裝飾，而是結構性的。

\subsection{Question \& Answer}

一句話要怎樣才能不只抓住讀者一下，而是真正留在他身上？

「抓住」只是短暫攫住眼睛；「餘響」則會在那一刻過後仍然留下來。依 Vuong 的說法，一句話之所以能留下，是因為它的句法形狀足以承受那份知覺壓力。這一行不能只是傳達概念，它還必須帶著節奏、延遲與壓力進入讀者。因此，抽象之物要變得可記，只有在它被賦予了身體之後。

\section{詩與自然書寫，作為語言的實驗室}

從讀者那種事後仍然被留下的感覺出發，講座再次擴大，這一次走向藝術養成。問題變成：句子究竟還能在哪裡冒險？Vuong 的回答是：詩與自然書寫，仍然保留了大多數當代散文已經部分放棄的自由。

第一條主張非常直接：
\begin{equation}
\text{指派被移除} \Longrightarrow \text{語言變得可供實驗}.
\end{equation}
第二條則順著它而來：
\begin{equation}
\text{詩}，\ \text{自然書寫} \Longrightarrow \text{句子在局部上擁有最高自由}.
\end{equation}
詩之所以是顯而易見的實驗室，是因為作者無須時時刻刻維持情節、說明性清晰或角色管理。自然書寫之所以成為一個平行實驗室，則出於更微妙的理由：在那裡，純粹的摹擬根本撐不住。我們早就看過草地，也早就知道泥巴看起來像什麼。那麼，若一句話只是把已經可見之物重複一次，我們為什麼還要讀它？

J.A. Baker 的段落便是對這個問題的現場回答。Baker 從平實的天氣與沼地描寫起手，接著卻進入一種不斷繁殖的泥巴句法，而那已經不再只是命名物質本身。泥變成了壓力、質地、危險、環境、心理場域。到了句尾，Baker 的內在性已經滲入物件本身：
\begin{equation}
\text{泥巴描寫} + \text{被釋放的內在性} \Longrightarrow \text{摹擬被打破}.
\end{equation}
這正是 Vuong 希望筆記能保留下來的運動。Baker 並不是離開世界，而是讓主觀力量改變了世界在頁面上的樣子。

主持人在這裡補上的另一個關鍵詞，是 fun。這個詞非常重要，否則整場講座就可能滑成一場全是理論與指責的演說。那盒彩色蠟筆裡頭有許多種藍色的快感，便成了一種藝術知覺的模型。當我們增加顏色層次、質地與用途時，我們不只是在做更細的分類，我們也是在重新打開驚奇。這也就是為什麼講座能夠毫不費力地從 Baker 的泥，走回兒童式的注意力。

Shklovsky 與 Tolstoy 的材料，也在這裡以另一個語域回來。Tolstoy 不說「白樺樹」，而是像在說：一棵頭髮蓬鬆的大樹，帶著發亮的白色樹幹。重點不在於描寫的冗長，而在於：要讓我們再次相信那個孩子。當命名太快封死物件時，定義本身就可能成為想像力的敵人。

也正因如此，講座後段 Eduardo C. Corral 那個關於 moss 的明喻才顯得如此重要。逐字稿在這個例子前後略有不穩，但其中技術主張卻非常清楚：這條比較之所以有效，不是因為 image 對 image，而是因為 behavior 對 behavior：
\begin{align}
\text{掌聲} \not\sim \text{苔蘚於圖像}, \qquad \text{掌聲} \sim \text{苔蘚於行為},\\
\text{行為對應} \Longrightarrow \text{苔蘚彷彿加速生長}.
\end{align}

\paragraph{Worked Example: Moss and Applause.}
\begin{enumerate}
\item 我們不再問苔蘚靜態地「像什麼」。
\item 我們改問：苔蘚如何動作，或者我們如何重新感知它的動作。
\item ``掌聲'' 把一種加快的、模糊的、集體性的運動借給了苔蘚。
\item 這種被轉移過來的行為感，使苔蘚彷彿在我們眼前移動起來。
\item 因此，這個明喻改變的不是目標的表面，而是目標被經驗到的時間性。
\end{enumerate}

Corral 花了九年才寫出一本短書的軼事，在這裡之所以重要，是因為它把從 Babel 開始打開的那個環再次關起來。強而有力的隱喻，從來都不是瞬間聰明，而是長久觀看之後的產物。

\subsection{Question \& Answer}

為什麼詩與自然書寫，比大多數當代散文更能為句子保留自由？

因為兩者都減輕了「指派」的即時壓力。詩可以讓語言本身成為焦點。自然書寫則因為無法只靠報告而存活——可見之物我們原本就知道——而逼迫句子朝向創造、陌生化與遊戲。所謂實驗室，不過就是一個作者被允許去發現世界不只一種用途的地方。

\section{出版時間、讀者時間，以及敢於不服從的倫理}

講座後段，便從局部工藝轉向更大的系統。期刊內部風格、教學上的犬儒、好萊塢式保守、那個「中西部讀者」的想像，以及 Lotman 的雙時間模型，全部都開始彙聚成同一個問題：為什麼文化在原則上讚美原創，卻又在實際操作裡用同質化來獎勵作者？

Vuong 早先已經說過句子如何被標準化。現在，他更進一步指出：制度是如何積極維持這種標準化的。The New Yorker 的例子之所以重要，不是因為它是單純的譴責。House style 當然教會人清楚，也同時製造出一種品牌辨識度。作者可能在那裡發表作品，卻很快發現：文章的思想仍然是自己的，但節奏卻不再完全屬於自己。重點不是說 house style 是邪惡的，而是說：清晰、有效與品牌期待，從來都不是中性的。

接著，那個關於「中西部讀者」的軼事，把這種犬儒推得更尖。一位編輯問：那樣的讀者真的能處理繁複、巴洛克式的散文嗎？Vuong 從這句話裡聽見的不是民主式的關懷，而是居高臨下的輕蔑。他說，普通讀者一樣有神經系統、有記憶，也有自己的閱讀生命。這個軼事，於是正好為 Uri Lotman 鋪路：
\begin{equation}
\text{閱讀} = (\text{共時線},\text{歷時線}).
\end{equation}
所謂共時閱讀，是嵌在某一季節、某一宣傳循環、某一市場年度、某一制度時刻之內的閱讀；歷時閱讀，則是與個人長年累積起來的閱讀經驗一起發生的閱讀。

這便產生了講座裡最有力的一條制度對比：
\begin{equation}
\text{系統以共時方式判斷}, \qquad \text{讀者以歷時方式判斷}.
\end{equation}
出版運作於季度之中，評論文化運作於季度之中，批評也常常運作於季度之中；但讀者不是。讀者也許在這本當季得獎小說到來之前，早就讀過 Shakespeare、Melville、Baldwin、Annie Dillard 或 Chaucer。因此，一本為了共時成功而被標準化的書，一旦落到歷時閱讀者手中，所面對的測試就完全不同。

主持人接著用 Rotten Tomatoes 的類比，把 Lotman 的理論翻譯成一個熟悉的文化分裂。評論家被嵌進了制度性的當下；觀眾則帶著較不受管控的比較累積而來。講座的意思並不是說觀眾純潔、評論家腐敗，而是說，這兩者在不同的時鐘上運作。

然後，Vuong 又從出版時間回到語言時間。定義有時會收縮事物，但也有時會擴張事物。好的字典與詞源學，會重新打開詞，而不是把它們封死。Passion 的例子在這裡極為決定性：今日語用把它聽成熱情或強度，但它其實攜帶著更古老的 suffering 歷史。一旦這段歷史被重新接上，那個詞就改變了。這也就是通向 Wittgenstein 的橋：
\begin{equation}
\text{詞語的意義} = \text{它的用法}.
\end{equation}
用法會改變定義，字典只是追著我們跑而已。這個教訓之所以對學生尤其重要，是因為對官方規則的恐懼，正是將寫作推回標準化的另一條路。

接下來，講座再次把範圍推大，從詞彙走向文化：
\begin{equation}
\text{邊緣的創新} \Longrightarrow \text{被中心捕捉} \Longrightarrow \text{被同質化的產品}.
\end{equation}
邊緣是不斷移動的，而中心則總是在捕捉它。Lotman 的同心圓文化模型解釋了：為什麼創新往往從被承認中心之外的地方開始，最後又作為被壓平的商品返回。這對詞語、風格、場景與形式都同樣成立。

接著，主持人問出了整場後段最實際的一個問題：那麼，作者到底該怎麼做？Vuong 的回答不是另一條技巧性規則，而是一種倫理：daringness 與 disobedience。敢，是願意賭、願意冒著失敗風險；不服從，則是拒絕只因為「回到隊伍裡」會被誇獎，就真的退回去。那個滑板的類比之所以重要，是因為它替 daringness 提供了一種身體感：人把自己朝一個動作丟出去，卻並不確定是否會成功落地。失敗不是意外副產品，而本來就屬於那套練習本身。

講座最後甚至越過工藝，走向哲學性的風險。語言確實塑造了 Vuong 的人生，它帶來謀生、帶來進入世界的權利，也帶來擴張知覺的能力；但同樣的語言也可能被暴政捕獲。報紙與廣播，往往正是威權首先要控制的東西。於是，最後這個悖論必須保持完整：
\begin{align}
\text{語言} &\Longrightarrow \text{權力} + \text{生計} + \text{知覺的擴張},\\
\text{語言} &\Longrightarrow \text{捕獲} + \text{暴政} + \text{同質化}.
\end{align}
Thomas Thistlewood 的圖書館、納粹軍官那種被培養出來的殘酷，以及其他歷史例子，都使得這場講座不至於滑向感傷。與此同時，Harriet Beecher Stowe 依然站在天平的另一端：文學有時確實會改變歷史。Vuong 堅持要同時保住這兩個真理。

\subsection{Question \& Answer}

為什麼文化一邊要求原創，一邊又獎勵順從？而作者究竟應該怎麼辦？

因為制度需要的是可辨認的產品。標準化更容易編輯、品牌化、流通，也更容易被稱讚。而作者的回答，在這場講座裡，不是某種純潔幻想，而是一套實踐：
\begin{enumerate}
\item 保護那個最初的驚奇矩陣，
\item 用歷時性的方式閱讀，而不只用季度性的方式閱讀，
\item 相信活生生的使用，而不是被官方規則嚇倒，
\item 同時鍛鍊敢與不服從，
\item 承認語言有力量，卻不假裝它是無辜的。
\end{enumerate}
講座並沒有保證這條路一定會被回報。它只提出一個更強的判斷：若低於這個標準，那句子就會膽怯到根本無法活下去。

\section{Summary}

這場講座是以穩定擴大的方式展開的。關於隱喻的問題，變成了關於工作坊的問題；關於工作坊的問題，又變成了關於句子歷史的問題；句子的歷史則打開了對陌生化、門檻、句法、模式、詩、自然書寫、出版時間、讀者時間，最後甚至是語言自身道德曖昧性的討論。而在每一步擴張之後，論證又都會再次縮回到一個局部例子裡：Babel 的落日、植物學家的方法、被移位的玫瑰、Siken 的星星、Baker 的泥、Corral 的苔蘚，以及那位被想像出來的「中西部讀者」。

最後留下來的，並不是一套鼓吹風格極大化的教條，而是一條更精確的注意力標準。我們被要求看得更久，懷疑那些僅僅可辨認的東西，聽見句子作為意識承載者的功能，並記住：形式上的大膽，既是一場藝術賭注，也是一場倫理賭注。這場講座理應在這裡結束：不是在解決中，而是在對語言究竟能做什麼、又不能保證什麼，抱持一種更高強度的嚴肅之中。